% ===========================================================================
\subsection{\file{new-design-reference/index.html}}
% ===========================================================================

\textbf{HTML, 597 lines.} The page itself: every word a visitor reads, and the
structure that holds those words. If you delete every other file in this
repository, this one still opens in a browser and still says everything the
project wants to say. It would simply look plain.

\begin{provenance}{new-design-reference/index.html}
Hand-written. It arrived complete in the very first commit, \code{0d77d51} on
2026-07-28 (``Launch ai.digitalise.agency''), at the path \file{index.html}, and
was moved to its present path in \code{da508a7} on 2026-07-29 with no change to
its contents. There is no generator banner, no framework scaffold, and no
\file{package.json} that could have produced it.

Its icons are a different matter and are worth naming precisely. Every icon in
the file is an inline SVG drawing whose path data matches the Heroicons outline
set, a widely used free icon library: the shapes, the \code{24 24} view box, the
\code{stroke-linecap} and \code{stroke-linejoin} attributes and the
\code{1.5}-based coordinates are that library's house conventions. \textbf{The
file carries no attribution comment saying so}, so this is stated as an inference
from the path data, not as a documented fact. \textbf{Confidence: high for
hand-written, inferred for the icon origin.}
\end{provenance}

\subsubsection*{Why it exists}

Without this file there is no page. The stylesheet styles class names that would
belong to nothing, and the JavaScript would search for elements that do not
exist and quietly do nothing, which is exactly how it is written to behave.

\subsubsection*{Line by line}

\paragraph{Lines 1 to 17: the head, which the visitor never sees.}

\begin{lstlisting}[language=HTML,firstnumber=1,caption={\file{index.html}, lines 1--17}]
<!DOCTYPE html>
<html lang="en">
<head>
<meta charset="UTF-8" />
<meta name="viewport" content="width=device-width, initial-scale=1.0" />
<title>Digitalise AI: Done-for-you AI systems for entrepreneurs</title>
<meta name="description" content="Done-for-you AI systems that answer, follow up, and publish, so your business keeps growing without adding headcount." />
<meta property="og:title" content="Digitalise AI" />
<meta property="og:description" content="Done-for-you AI systems that answer, follow up, and publish." />
<meta property="og:type" content="website" />
<meta name="theme-color" content="#05070d" />

<link rel="preconnect" href="https://fonts.googleapis.com" />
<link rel="preconnect" href="https://fonts.gstatic.com" crossorigin />
<link href="https://fonts.googleapis.com/css2?family=Space+Grotesk:wght@500;600;700&family=Inter:wght@400;500;600;700&display=swap" rel="stylesheet" />
<link rel="stylesheet" href="/css/site.css" />
</head>
\end{lstlisting}

Line 1 is not a tag at all. \code{<!DOCTYPE html>} is a declaration that tells the
browser ``interpret this as modern HTML''. It is a historical artefact: without
it browsers fall back to a compatibility mode that imitates the bugs of software
from the 1990s. It must be the first thing in the file, and it takes no closing
partner.

Line 2 opens the document and declares its language as English with
\code{lang="en"}. That attribute is small and does real work: screen readers use
it to choose a pronunciation, and search engines use it to decide who to show the
page to.

Line 3 opens the \code{<head>}. Everything between here and line 17 is
information \emph{about} the page rather than content \emph{of} it. Nothing in
this block appears on screen.

Line 4 declares the character encoding as UTF-8, which is the scheme that lets a
file contain any character in any language. This matters concretely here: the
page contains a middle dot and a heart symbol, and without this line those would
render as mojibake. It should appear as early as possible, because the browser
has to guess the encoding until it is told.

\begin{jargon}{Meta tag}
A tag inside \code{<head>} carrying a fact about the page rather than content for
the page. It has no closing partner and displays nothing. Different meta tags are
read by completely different audiences: browsers, search engines, and social
networks each look for their own.
\end{jargon}

Line 5 is the viewport declaration and it is the single line that makes the page
work on a phone. \code{width=device-width} tells the browser to treat the page
width as the real width of the device, and \code{initial-scale=1.0} says do not
zoom out to start with. Without it, a phone pretends to be about 980 pixels wide
and shrinks everything, which is why old sites look like tiny unreadable
newspapers on a phone. Every media query in the stylesheet depends on this line
being present.

Line 6 is the \code{<title>}: the text in the browser tab and the blue headline
in a search result. Note it for later, because it is how section 1.2 proved this
page is not the live one.

Line 7 is the description, which search engines often show as the grey text under
the title. It is 118 characters, comfortably inside the roughly 155 that get
displayed.

Lines 8 to 10 are Open Graph tags, a convention invented by Facebook and now read
by nearly every chat app and social network to build the preview card when
somebody pastes a link. Note that \code{og:title} here says only ``Digitalise
AI'', shorter than the real title on line 6, which is a sensible choice: preview
cards truncate.

\begin{honest}{The Open Graph set is incomplete in the way that matters most}
The three tags present are \code{og:title}, \code{og:description} and
\code{og:type}. The two missing ones are \code{og:image} and \code{og:url}.
Without \code{og:image}, a pasted link produces a preview card with no picture,
which is the difference between a link that gets clicked and one that does not.
\file{ASSETS.md} lists \code{og-image.png} as required before launch, and it was
never created. See section 7.5.
\end{honest}

Line 11 sets \code{theme-color} to \code{\#05070d}, the same near-black used as
the page background. Mobile browsers paint their address bar this colour, so the
browser chrome blends into the page instead of framing it in white.

Line 12 is blank. Blank lines have no effect on how HTML renders; they are purely
for the human reading the file, and here the author uses them to separate the
page's identity from its external resources.

Lines 13 and 14 are \code{preconnect} hints. They tell the browser: you will
shortly need to talk to these two Google servers, so start the handshake now
rather than waiting. Establishing a secure connection takes several network round
trips, and doing that early rather than at the moment of need typically saves a
few hundred milliseconds. The \code{crossorigin} attribute on line 14 is required
on the second one, because font files are fetched under stricter rules than
stylesheets, and a preconnect made without it would open a connection the font
request cannot reuse. That subtlety catches out experienced developers
constantly.

Line 15 actually fetches the fonts: two typefaces, Space Grotesk in three weights
and Inter in four. \code{display=swap} at the end is the important part. It tells
the browser to render text immediately in a fallback typeface and swap in the
real one when it arrives, rather than hiding the text for up to three seconds
while it waits. The visible cost is that text reflows slightly; the benefit is
that a visitor on a slow connection can read the page at all.

\begin{jargon}{Font weight}
How heavy a typeface looks, on a numeric scale where 400 is normal and 700 is
bold. Requesting \code{wght@500;600;700} downloads three separate files. Each
weight requested is another download, which is why the list is short.
\end{jargon}

Line 16 links the stylesheet. This is the path discussed in section 11: the
leading slash in \code{/css/site.css} means ``from the top of the site'', which
is correct when served by a web server and wrong when the file is opened directly
off a disk.

Line 17 closes the head.

\paragraph{Lines 18 to 35: the navigation bar.}

\begin{lstlisting}[language=HTML,firstnumber=18,caption={\file{index.html}, lines 18--35}]
<body>

<!-- ══════════════════════════════════════════════
     NAV
══════════════════════════════════════════════ -->
<nav class="nav" aria-label="Site navigation">
  <div class="nav-inner">
    <a href="/" class="nav-logo" aria-label="Digitalise AI home">
      <svg viewBox="0 0 28 28" fill="none" aria-hidden="true">
        <rect width="28" height="28" rx="8" fill="#1967ff"/>
        <path d="M7 14c0-3.866 3.134-7 7-7s7 3.134 7 7-3.134 7-7 7" stroke="#fff" stroke-width="2.2" stroke-linecap="round"/>
        <circle cx="14" cy="14" r="2.5" fill="#fff"/>
      </svg>
      digitalise <span class="logo-ai">AI</span>
    </a>
    <a href="#book" class="btn btn-primary">Book a call</a>
  </div>
</nav>
\end{lstlisting}

Line 18 opens \code{<body>}: everything from here to line 596 is what a visitor
actually sees.

Line 19 is blank, and lines 20 to 22 are a comment banner. Anything between
\code{<!-{}-} and \code{-{}->} is ignored entirely by the browser. This file uses
seven of these banners as signposts, one per major band of the page, each drawn
with a rule of repeated double-line characters. They are the author's table of
contents, and they are the reason a 597-line file stays navigable.

Line 23 opens the navigation bar. Two things are happening. \code{<nav>} is a
\emph{semantic} tag: it does not look different from a plain \code{<div>}, but it
announces to screen readers and search engines that this is the site's
navigation. \code{aria-label="Site navigation"} gives that region a spoken name,
so a screen-reader user hearing a list of page regions hears something meaningful.

\begin{jargon}{ARIA}
Accessible Rich Internet Applications: a set of attributes, all starting with
\code{aria-}, that describe a page to assistive software such as screen readers.
They change nothing visually. \code{aria-label} supplies a name;
\code{aria-hidden="true"} tells assistive software to skip something entirely
because it is purely decorative.
\end{jargon}

Line 24 opens a \code{<div>} labelled \code{nav-inner}. A \code{<div>} is a
generic container with no meaning at all; it exists purely to be styled. The
reason for a second container inside \code{<nav>} is a common layout pattern: the
outer element spans the full window width so its background can stretch edge to
edge, while the inner one is capped at a maximum width and centred, so the logo
and button line up with the content further down the page.

Line 25 opens a link to \code{/}, the site's home page. Since this \emph{is} the
home page it reloads the page, which is the conventional behaviour of a logo.
Its \code{aria-label} matters here: the link's visible content is a drawing plus
the words ``digitalise AI'', and without the label a screen reader would announce
the drawing awkwardly.

Lines 26 to 30 are the logo, drawn in SVG.

\begin{jargon}{SVG}
Scalable Vector Graphics: a picture described as shapes and coordinates rather
than as a grid of coloured dots. Because it is instructions rather than pixels,
it stays perfectly sharp at any size, and because it is text it can sit directly
inside the HTML file with no extra download.
\end{jargon}

\code{viewBox="0 0 28 28"} sets up a private coordinate system 28 units square;
every coordinate inside is relative to that, so the whole drawing can be scaled to
any size by CSS without touching the numbers. \code{fill="none"} sets the default
to no fill, which each shape then overrides. \code{aria-hidden="true"} hides it
from screen readers, correctly: the link already has a text label, and announcing
the drawing too would be duplication.

Line 27 draws a 28 by 28 rounded square (\code{rx="8"} is the corner radius) in
the brand blue. Line 28 draws the swoosh: \code{d} is a path description in a
compact language where \code{M7 14} means ``move to the point (7,14)'' and the
\code{c} and \code{s} commands that follow are curve instructions. It is written
to be generated, not read, and nobody edits it by hand. Line 29 puts a white dot
in the middle. Line 30 closes the drawing.

Line 31 is the wordmark. ``digitalise'' is plain text and ``AI'' is wrapped in a
\code{<span>} labelled \code{logo-ai} so the stylesheet can colour those two
letters blue. A \code{<span>} is the inline equivalent of a \code{<div>}: no
meaning, no visual effect, purely a hook for styling.

Line 32 closes the logo link. Line 33 is the navigation's only button, a link to
\code{\#book}. Note it is a link and not a \code{<button>}, which is correct: it
navigates somewhere rather than performing an action. It carries two classes,
\code{btn} for the shared button shape and \code{btn-primary} for the filled blue
variant. That two-class pattern repeats throughout the file and is explained in
section 7.2.

Lines 34 and 35 close the containers.

\paragraph{Lines 36 to 57: the hero opens, and the first floating tile.}

\begin{lstlisting}[language=HTML,firstnumber=36,caption={\file{index.html}, lines 36--57}]


<!-- ══════════════════════════════════════════════
     HERO
══════════════════════════════════════════════ -->
<section class="hero" aria-label="Hero">

  <!-- animated background tiles -->
  <div class="hero-bg" aria-hidden="true">

    <!-- Tile 1: LinkedIn Video -->
    <div class="tile tile-blue">
      <div class="tile-thumb">
        <svg viewBox="0 0 24 24" fill="none" stroke="currentColor" stroke-width="1.8" aria-hidden="true">
          <path stroke-linecap="round" stroke-linejoin="round" d="M15.75 10.5l4.72-4.72a.75.75 0 011.28.53v11.38a.75.75 0 01-1.28.53l-4.72-4.72M4.5 18.75h9a2.25 2.25 0 002.25-2.25v-9a2.25 2.25 0 00-2.25-2.25h-9A2.25 2.25 0 002.25 7.5v9a2.25 2.25 0 002.25 2.25z"/>
        </svg>
      </div>
      <div class="tile-meta">
        <span class="tile-label">LinkedIn Video</span>
        <span class="tile-metric">847K views</span>
      </div>
    </div>
\end{lstlisting}

Lines 36 to 39 are a blank line and the second comment banner.

Line 40 opens the hero, the full-screen band at the top of the page.
\code{<section>} is another semantic tag, meaning ``a distinct part of the
document''.

Line 42 is a comment, and line 43 opens the container for the eight drifting
cards. \code{aria-hidden="true"} on this container is the single most important
accessibility decision in the file: it hides all eight tiles, and every number
inside them, from screen readers. That is correct, because they are decoration.
It is also the reason the invented numbers inside them are a presentation problem
rather than an accessibility one.

Lines 46 to 57 are the first tile. \code{class="tile tile-blue"} again uses the
two-class pattern: \code{tile} carries the shared card shape and the floating
animation, \code{tile-blue} adds only the blue gradient behind the icon.

Inside, line 47 opens the thumbnail area, and lines 48 to 51 draw a video-camera
icon. \code{stroke="currentColor"} on line 48 is worth understanding: it means
``draw the outline in whatever text colour applies here'', so the icon inherits
its colour from CSS instead of hard-coding one. That single attribute is why one
icon definition can appear in white on one background and blue on another.

Lines 53 to 56 are the tile's caption row: a label on the left and a metric on
the right.

\begin{honest}{The tile metrics are invented}
\code{847K views} is not a measurement. Neither are the other five numbers in
these tiles: 234K impressions, 1.2M views, 512K views, 1.8M views and 19K
reactions. Nothing in this repository or its history connects them to a real
account or a real campaign, and they appeared fully formed in the first commit
alongside the rest of the design.

They are decorative and hidden from assistive software, and a designer would call
them realistic placeholder content, which is a normal thing to put in a mockup.
But a sighted visitor cannot tell decoration from a claim, and reads them as
results this agency has produced. That is worth raising specifically because the
agency's main website was built to a stated standard of using no fabricated
numbers.
\end{honest}

\paragraph{Lines 58 to 70: the second tile, built from bars instead of an icon.}

\begin{lstlisting}[language=HTML,firstnumber=58,caption={\file{index.html}, lines 58--70}]

    <!-- Tile 2: Post tile -->
    <div class="tile tile-orange">
      <div class="tile-lines" style="padding:10px 0;">
        <div class="tile-line"></div>
        <div class="tile-line"></div>
        <div class="tile-line"></div>
      </div>
      <div class="tile-meta">
        <span class="tile-label">LinkedIn Post</span>
        <span class="tile-metric">234K impr.</span>
      </div>
    </div>
\end{lstlisting}

This tile represents a text post rather than a video, so instead of an icon it
shows three grey bars standing in for lines of text. Each \code{<div
class="tile-line">} is completely empty: it has no content at all, and exists
only so that CSS can give it a height, a width and a rounded grey background.
This is a standard technique for drawing an abstract shape without an image, and
the widths (100\%, 85\%, 60\%) are set in the stylesheet at lines 216 to 218 to
suggest a paragraph ending mid-line.

\begin{honest}{An inline style that fights the stylesheet}
\code{style="padding:10px 0;"} on line 61 sets the spacing directly on the
element. The stylesheet also has a rule for this class, \code{.tile-lines
\{ padding: 0; \}} at line 209, and an inline style always wins over a stylesheet
rule. So that rule never has any effect, in either of the two places the class is
used. One of the two should go; the stylesheet rule is the dead one.
\end{honest}

\begin{jargon}{Inline style}
Appearance written directly on an element with a \code{style} attribute, rather
than in the stylesheet. It overrides almost anything the stylesheet says, which
makes it useful for one-offs and dangerous as a habit: rules written in one file
stop being able to control what appears in the other.
\end{jargon}

\paragraph{Lines 71 to 80: the third tile, a statistic.}

\begin{lstlisting}[language=HTML,firstnumber=71,caption={\file{index.html}, lines 71--80}]

    <!-- Tile 3: Stats tile -->
    <div class="tile" style="background:var(--s1);">
      <div class="tile-stat-num">1.2M</div>
      <div class="tile-stat-lbl">Views this month</div>
      <div class="tile-meta" style="margin-top:12px;">
        <span class="tile-label">AI Campaign</span>
        <span class="tile-metric">+340%</span>
      </div>
    </div>
\end{lstlisting}

This tile has no icon and no bars, just a large number, a caption and the usual
meta row. Note \code{style="background:var(-{}-s1);"} on line 73: an inline style
that reads a custom property defined in the stylesheet. It works, and it is
redundant, because the \code{.tile} rule at line 155 already sets exactly that
background. Two more tiles (lines 96 and 121) repeat the same redundant inline
background.

\paragraph{Lines 81 to 93: the fourth tile.}

\begin{lstlisting}[language=HTML,firstnumber=81,caption={\file{index.html}, lines 81--93}]

    <!-- Tile 4: Video -->
    <div class="tile tile-teal">
      <div class="tile-thumb">
        <svg viewBox="0 0 24 24" fill="none" stroke="currentColor" stroke-width="1.8" aria-hidden="true">
          <path stroke-linecap="round" stroke-linejoin="round" d="M5.25 5.653c0-.856.917-1.398 1.667-.986l11.54 6.348a1.125 1.125 0 010 1.971l-11.54 6.347a1.125 1.125 0 01-1.667-.985V5.653z"/>
        </svg>
      </div>
      <div class="tile-meta">
        <span class="tile-label">AI Video</span>
        <span class="tile-metric">512K views</span>
      </div>
    </div>
\end{lstlisting}

Structurally identical to tile 1, with two differences: the colour class is
\code{tile-teal}, and the icon is a play triangle rather than a video camera.
That same play-triangle path appears six more times in this file, once on each
portfolio card. Nothing shares it: the same 180 characters of path data are
repeated seven times in total, which is the cost of having no templating system.
It is a real duplication, and it is also the correct trade for a project that
deliberately has no build step.

\paragraph{Lines 94 to 105: the fifth tile, a status badge.}

\begin{lstlisting}[language=HTML,firstnumber=94,caption={\file{index.html}, lines 94--105}]

    <!-- Tile 5: Published badge -->
    <div class="tile" style="background:var(--s1);">
      <div style="display:flex;align-items:center;gap:8px;margin-bottom:10px;">
        <svg viewBox="0 0 24 24" fill="none" stroke="#34d399" stroke-width="2" width="20" height="20" aria-hidden="true">
          <path stroke-linecap="round" stroke-linejoin="round" d="M9 12.75L11.25 15 15 9.75M21 12a9 9 0 11-18 0 9 9 0 0118 0z"/>
        </svg>
        <span style="font-family:var(--font-d);font-size:13px;font-weight:700;color:#34d399;">Published</span>
      </div>
      <div class="tile-stat-num" style="font-size:22px;">47 posts</div>
      <div class="tile-stat-lbl">this week, zero manual work</div>
    </div>
\end{lstlisting}

This is the most heavily inline-styled element in the file: line 97 lays out a
row, line 101 styles the word ``Published'', and line 103 overrides a font size.
None of it has a class, so none of it can be restyled from the stylesheet.

Note the green on lines 98 and 101 is written as the literal \code{\#34d399},
even though the stylesheet defines exactly that colour as \code{-{}-green} at line
15. Change the token and this tile would not follow. The same literal appears
again at line 276 in the pipeline. This is the clearest illustration in the file
of what inline styles cost: a colour system that only mostly applies.

\paragraph{Lines 106 to 118: the sixth tile.}

\begin{lstlisting}[language=HTML,firstnumber=106,caption={\file{index.html}, lines 106--118}]

    <!-- Tile 6: Violet video -->
    <div class="tile tile-violet">
      <div class="tile-thumb">
        <svg viewBox="0 0 24 24" fill="none" stroke="currentColor" stroke-width="1.8" aria-hidden="true">
          <path stroke-linecap="round" stroke-linejoin="round" d="M15.75 10.5l4.72-4.72a.75.75 0 011.28.53v11.38a.75.75 0 01-1.28.53l-4.72-4.72M4.5 18.75h9a2.25 2.25 0 002.25-2.25v-9a2.25 2.25 0 00-2.25-2.25h-9A2.25 2.25 0 002.25 7.5v9a2.25 2.25 0 002.25 2.25z"/>
        </svg>
      </div>
      <div class="tile-meta">
        <span class="tile-label">AI Video</span>
        <span class="tile-metric">1.8M views</span>
      </div>
    </div>
\end{lstlisting}

Identical in structure to tile 1, and using the identical video-camera icon path,
with a violet gradient and a different number. This is the third of four tiles
built from the same template.

\paragraph{Lines 119 to 130: the seventh tile, a schedule.}

\begin{lstlisting}[language=HTML,firstnumber=119,caption={\file{index.html}, lines 119--130}]

    <!-- Tile 7: Schedule -->
    <div class="tile" style="background:var(--s1);">
      <div style="display:flex;align-items:center;gap:8px;margin-bottom:10px;">
        <svg viewBox="0 0 24 24" fill="none" stroke="var(--accentH)" stroke-width="1.8" width="18" height="18" aria-hidden="true">
          <path stroke-linecap="round" stroke-linejoin="round" d="M6.75 3v2.25M17.25 3v2.25M3 18.75V7.5a2.25 2.25 0 012.25-2.25h13.5A2.25 2.25 0 0121 7.5v11.25m-18 0A2.25 2.25 0 005.25 21h13.5A2.25 2.25 0 0021 18.75m-18 0v-7.5A2.25 2.25 0 015.25 9h13.5A2.25 2.25 0 0121 11.25v7.5"/>
        </svg>
        <span style="font-family:var(--font-d);font-size:12px;font-weight:600;color:var(--accentH);">Scheduled</span>
      </div>
      <div class="tile-stat-num" style="font-size:22px;">34 posts</div>
      <div class="tile-stat-lbl">queued this week</div>
    </div>
\end{lstlisting}

The mirror image of tile 5: same layout, calendar icon instead of a tick, and
``queued'' instead of ``published''. Together the pair tells a small story in the
background of the hero, that work is both finished and lined up. Note that here
the colour \emph{is} read from a token, \code{var(-{}-accentH)} on lines 123 and
126, unlike tile 5 which used a literal. The two tiles were written minutes apart
and are inconsistent with each other.

\paragraph{Lines 131 to 145: the eighth tile, and the end of the backdrop.}

\begin{lstlisting}[language=HTML,firstnumber=131,caption={\file{index.html}, lines 131--145}]

    <!-- Tile 8: Reactions -->
    <div class="tile tile-blue">
      <div class="tile-lines" style="padding:10px 0;">
        <div class="tile-line"></div>
        <div class="tile-line"></div>
        <div class="tile-line" style="width:50%;"></div>
      </div>
      <div class="tile-meta">
        <span class="tile-label">Viral Post</span>
        <span class="tile-metric">19K ❤</span>
      </div>
    </div>

  </div><!-- /hero-bg -->
\end{lstlisting}

The last tile repeats the bar pattern from tile 2, overriding the third bar's
width to 50\% inline so the two text tiles do not look identical. Line 141 is
where the heart character appears; as section 1.4 explained, this document prints
it as a heart symbol.

Line 145 closes the tile container, with a trailing comment naming what was
closed. That habit is worth copying: in a file with this much nesting, a bare
\code{</div>} tells the reader nothing.

\begin{keyidea}{Why the tile order in the file is not the order on screen}
The eight tiles are written top to bottom in the file, but they appear scattered
around the hero. Nothing here places them. Each one is positioned by the
stylesheet, using rules that count children: \code{.tile:nth-child(1)},
\code{(2)} and so on, at lines 162 to 169, each setting a corner and a delay
before that tile starts drifting. \textbf{This means the order of these blocks in
the HTML is load-bearing.} Move tile 5 above tile 3 in this file, without touching
the CSS, and both tiles jump to new positions on screen and start animating at
different times.
\end{keyidea}

\paragraph{Lines 146 to 182: the overlay and the hero's actual content.}

\begin{lstlisting}[language=HTML,firstnumber=146,caption={\file{index.html}, lines 146--182}]

  <div class="hero-overlay" aria-hidden="true"></div>

  <!-- hero content -->
  <div class="hero-content">
    <div class="hero-badge">
      <span class="hero-badge-dot"></span>
      Digitalise Agency
    </div>

    <h1 class="hero-h1">
      AI systems built<br>
      to grow while you <em>sleep.</em>
    </h1>

    <div class="hero-actions">
      <a href="#book" class="btn btn-primary">
        Book a call
        <svg width="16" height="16" viewBox="0 0 24 24" fill="none" stroke="currentColor" stroke-width="2.5" aria-hidden="true">
          <path stroke-linecap="round" stroke-linejoin="round" d="M13.5 4.5L21 12m0 0l-7.5 7.5M21 12H3"/>
        </svg>
      </a>
      <a href="#portfolio" class="btn btn-ghost">See the work</a>
    </div>

    <div class="hero-social-proof">
      <div class="avatars" aria-hidden="true">
        <div class="avatar av1">K</div>
        <div class="avatar av2">H</div>
        <div class="avatar av3">M</div>
        <div class="avatar av4">A</div>
      </div>
      <p class="hero-proof-text"><strong>17+ businesses</strong> growing with Digitalise AI</p>
    </div>
  </div>

</section>
\end{lstlisting}

Line 147 is an empty \code{<div>} that exists purely to darken the middle of the
hero so the headline stays readable over the drifting tiles. The stylesheet fills
it with two stacked gradients (lines 232 to 243). An empty element used as a
visual layer is a common and legitimate technique.

Lines 150 to 153 are the small pill above the headline. \code{hero-badge-dot} is
another empty element, styled into a small circle that pulses.

Lines 156 to 159 are the headline, in an \code{<h1>}. A page should have exactly
one \code{<h1>}, and this page does: it is the document's title as far as a screen
reader or a search engine is concerned. \code{<br>} on line 157 forces a line
break so the headline splits where the designer wants rather than where the window
happens to be narrow.

\begin{honest}{\code{<em>} is used for colour, not for emphasis}
Line 158 wraps ``sleep.'' in \code{<em>}, and the stylesheet then removes the
italics that tag normally implies and colours the word blue instead (lines 285 to
288). \code{<em>} means emphasis, and a screen reader may change its intonation
for it. Using it as a colour hook works, and it slightly misrepresents the
content to assistive software; a \code{<span>} with a class would have been the
honest choice. This is a small point and the intent is plainly visual.
\end{honest}

Lines 161 to 169 hold the two hero buttons. The first is the primary blue one,
containing both text and an arrow icon; it points at \code{\#book}. The second is
the ghost variant, pointing at \code{\#portfolio}. Both are ordinary links, and
both are intercepted by the JavaScript at line 70 to scroll smoothly rather than
jump.

Lines 171 to 179 are the social proof row: four overlapping circles bearing the
letters K, H, M and A, and a line of text. The circles are \code{aria-hidden}
because they carry no information.

\begin{honest}{The four avatar initials do not correspond to anything}
K, H, M and A do not match the three clients named further down the page (Hi-Tec
Bearings, DigiMarvixx, Synergy Event, whose own avatars use K, D and S). They are
placeholder initials that were never replaced. The ``17+ businesses'' claim beside
them matches the \code{data-target="17"} in the metrics band, so that number at
least is internally consistent; whether it is externally true cannot be
established from this repository.
\end{honest}

Line 182 closes the hero section.

\paragraph{Lines 183 to 210: the metrics band.}

\begin{lstlisting}[language=HTML,firstnumber=183,caption={\file{index.html}, lines 183--210}]


<!-- ══════════════════════════════════════════════
     METRICS BAND
     UPDATE: replace data-target values with real numbers
══════════════════════════════════════════════ -->
<section class="metrics" aria-label="Results">
  <div class="container">
    <div class="metrics-grid stagger">
      <div class="metric-item" style="--i:0">
        <div class="metric-num" data-target="17" data-suffix="+">0</div>
        <div class="metric-label">Projects Delivered</div>
      </div>
      <div class="metric-item" style="--i:1">
        <div class="metric-num" data-target="1.6" data-suffix="M+">0</div>
        <div class="metric-label">Organic Views for Clients</div>
      </div>
      <div class="metric-item" style="--i:2">
        <div class="metric-num" data-target="10" data-suffix="K+">0</div>
        <div class="metric-label">Users Served</div>
      </div>
      <div class="metric-item" style="--i:3">
        <div class="metric-num" data-target="300" data-suffix="%">0</div>
        <div class="metric-label">Avg. Traffic Growth</div>
      </div>
    </div>
  </div>
</section>
\end{lstlisting}

This is the first band that the JavaScript genuinely drives, and it introduces
the two mechanisms the rest of the page reuses.

The first is the \textbf{data attribute}. Look at line 193: the element's visible
content is the single character \code{0}, and the real number lives in
\code{data-target="17"}, with \code{data-suffix="+"} beside it.

\begin{jargon}{Data attribute}
Any attribute whose name starts with \code{data-}. HTML ignores it completely and
displays nothing; it exists so that JavaScript can attach its own information to
an element. \file{main.js} reads these two at lines 20 and 21 to know what number
to count towards and what to print after it.
\end{jargon}

Why is the visible content \code{0} rather than \code{17}? Because that is what a
visitor sees before the counter runs. It is also what a visitor sees
\emph{forever} if JavaScript is switched off or fails to load, which is the
weakness of this pattern: the page would show four zeros and claim nothing. The
robust alternative is to write the real number in the HTML and have the script
count towards it from zero, which fails gracefully. This file chose the fragile
form.

The second mechanism is the \textbf{stagger}. Line 191 gives the grid the class
\code{stagger}, and each child carries \code{style="-{}-i:0"}, \code{-{}-i:1} and so
on. That custom property is multiplied by 90 milliseconds in the stylesheet (line
90) to delay each child's fade-in, so the four items appear in sequence rather
than together. Setting a custom property inline is the cleanest way to pass a
per-element number to CSS, and it is used four times in this file.

Note that this band still carries an \code{UPDATE} comment on line 187 telling the
reader to replace the values with real numbers, even though commit \code{2c66a7a}
claimed to have done exactly that. See section 14.

\paragraph{Lines 211 to 235: the pipeline section and its first node.}

\begin{lstlisting}[language=HTML,firstnumber=211,caption={\file{index.html}, lines 211--235}]


<!-- ══════════════════════════════════════════════
     HOW IT WORKS (PIPELINE)
══════════════════════════════════════════════ -->
<section class="pipeline-section" aria-label="How it works">
  <div class="container">
    <div class="reveal">
      <p class="section-kicker">How it works</p>
      <h2 class="section-title">From idea to live<br>in one click.</h2>
    </div>

    <div class="pipeline" aria-label="Content production pipeline">

      <!-- Node 1 -->
      <div class="pipeline-node">
        <div class="pipeline-icon">
          <svg viewBox="0 0 24 24" fill="none" stroke="currentColor" stroke-width="1.8" aria-hidden="true">
            <path stroke-linecap="round" stroke-linejoin="round" d="M12 18v-5.25m0 0a6.01 6.01 0 001.5-.189m-1.5.189a6.01 6.01 0 01-1.5-.189m3.75 7.478a12.06 12.06 0 01-4.5 0m3.75 2.355a14.998 14.998 0 01-3 0M14.25 18v-.192c0-.983.658-1.823 1.508-2.316a7.5 7.5 0 10-7.517 0c.85.493 1.509 1.333 1.509 2.316V18"/>
          </svg>
        </div>
        <span class="pipeline-lbl">Your Idea</span>
      </div>

      <div class="pipeline-connector" aria-hidden="true"></div>
\end{lstlisting}

Lines 218 to 221 establish a pattern used by all four remaining sections: a
wrapper with the class \code{reveal}, containing a small uppercase ``kicker'' line
and a large heading. The \code{reveal} class is what the JavaScript watches for at
line 56, fading the pair upward into place when they scroll into view.

Line 220 is an \code{<h2>}. Heading levels form the document's outline: the single
\code{<h1>} in the hero, then an \code{<h2>} per major section. Getting that order
right is how a screen-reader user navigates a long page, and this file gets it
right throughout.

Line 223 opens the pipeline, and line 235 is the first connector: an entirely
empty element that the stylesheet turns into a short line between two nodes, and
then fills with a blue gradient when the section becomes visible. The nodes and
connectors alternate, which is why the stylesheet targets odd children for nodes
and even children for connectors (lines 400 to 426).

\paragraph{Lines 236 to 285: the remaining four pipeline nodes.}

\begin{lstlisting}[language=HTML,firstnumber=236,caption={\file{index.html}, lines 236--285}]

      <!-- Node 2 -->
      <div class="pipeline-node">
        <div class="pipeline-icon">
          <svg viewBox="0 0 24 24" fill="none" stroke="currentColor" stroke-width="1.8" aria-hidden="true">
            <path stroke-linecap="round" stroke-linejoin="round" d="M16.862 4.487l1.687-1.688a1.875 1.875 0 112.652 2.652L10.582 16.07a4.5 4.5 0 01-1.897 1.13L6 18l.8-2.685a4.5 4.5 0 011.13-1.897l8.932-8.931zm0 0L19.5 7.125M18 14v4.75A2.25 2.25 0 0115.75 21H5.25A2.25 2.25 0 013 18.75V8.25A2.25 2.25 0 015.25 6H10"/>
          </svg>
        </div>
        <span class="pipeline-lbl">AI Script</span>
      </div>

      <div class="pipeline-connector" aria-hidden="true"></div>

      <!-- Node 3 -->
      <div class="pipeline-node">
        <div class="pipeline-icon">
          <svg viewBox="0 0 24 24" fill="none" stroke="currentColor" stroke-width="1.8" aria-hidden="true">
            <path stroke-linecap="round" stroke-linejoin="round" d="M15.75 10.5l4.72-4.72a.75.75 0 011.28.53v11.38a.75.75 0 01-1.28.53l-4.72-4.72M4.5 18.75h9a2.25 2.25 0 002.25-2.25v-9a2.25 2.25 0 00-2.25-2.25h-9A2.25 2.25 0 002.25 7.5v9a2.25 2.25 0 002.25 2.25z"/>
          </svg>
        </div>
        <span class="pipeline-lbl">AI Video</span>
      </div>

      <div class="pipeline-connector" aria-hidden="true"></div>

      <!-- Node 4 -->
      <div class="pipeline-node">
        <div class="pipeline-icon">
          <svg viewBox="0 0 24 24" fill="none" stroke="currentColor" stroke-width="1.8" aria-hidden="true">
            <path stroke-linecap="round" stroke-linejoin="round" d="M6 12L3.269 3.126A59.768 59.768 0 0121.485 12 59.77 59.77 0 013.27 20.876L5.999 12zm0 0h7.5"/>
          </svg>
        </div>
        <span class="pipeline-lbl">Auto-Post</span>
      </div>

      <div class="pipeline-connector" aria-hidden="true"></div>

      <!-- Node 5 -->
      <div class="pipeline-node">
        <div class="pipeline-icon" style="background:rgba(52,211,153,.08);border-color:rgba(52,211,153,.25);">
          <svg viewBox="0 0 24 24" fill="none" stroke="#34d399" stroke-width="1.8" aria-hidden="true">
            <path stroke-linecap="round" stroke-linejoin="round" d="M9 12.75L11.25 15 15 9.75M21 12a9 9 0 11-18 0 9 9 0 0118 0z"/>
          </svg>
        </div>
        <span class="pipeline-lbl" style="color:#34d399;">Live</span>
      </div>

    </div><!-- /pipeline -->
  </div>
</section>
\end{lstlisting}

Four more nodes, each built to the identical shape: a container, an icon box, an
SVG, and a label. The icons in order are a lightbulb (an idea), a pencil on a
document (a script), a video camera (a video, reusing the same path as tile 1 and
tile 6), and a paper plane (posting).

Node 5 is the only one that differs, and it differs three times over: it overrides
its background and border colour inline on line 275, hard-codes the same green in
the SVG stroke on line 276, and colours its label inline on line 280. Everything
is written to make the final step green rather than blue, and none of it uses the
\code{-{}-green} token that exists for exactly this purpose. The same criticism as
tile 5 applies, and this is the fourth and last place the literal \code{\#34d399}
appears.

The five steps read as a claim about the product: an idea becomes a script,
becomes a video, gets posted automatically, and is live. That is the same pipeline
the Gantry project in this workspace actually implements, so the page is
advertising a real capability rather than an invented one, though nothing in this
repository connects the two.

\paragraph{Lines 286 to 318: the portfolio section and its first card.}

\begin{lstlisting}[language=HTML,firstnumber=286,caption={\file{index.html}, lines 286--318}]


<!-- ══════════════════════════════════════════════
     PORTFOLIO
     UPDATE: replace gradient .pt* classes with real video poster images.
     Add poster="path/to/thumb.jpg" to each <video> or swap divs for <img>.
══════════════════════════════════════════════ -->
<section class="portfolio-section" id="portfolio" aria-label="Work portfolio">
  <div class="container">
    <div class="reveal" style="margin-bottom:48px;">
      <p class="section-kicker">The work</p>
      <h2 class="section-title">Real content.<br>Real results.</h2>
    </div>

    <div class="portfolio-grid stagger">

      <div class="port-card" style="--i:0">
        <div class="port-thumb pt1">
          <div class="port-badge">847K views</div>
          <div class="port-thumb-overlay">
            <div class="port-play">
              <svg viewBox="0 0 24 24" fill="white" aria-hidden="true"><path d="M5.25 5.653c0-.856.917-1.398 1.667-.986l11.54 6.348a1.125 1.125 0 010 1.971l-11.54 6.347a1.125 1.125 0 01-1.667-.985V5.653z"/></svg>
            </div>
          </div>
        </div>
        <div class="port-info">
          <div class="port-platform">
            <svg viewBox="0 0 24 24" fill="none" stroke="currentColor" stroke-width="1.8" aria-hidden="true"><path stroke-linecap="round" stroke-linejoin="round" d="M13.5 4.5L21 12m0 0l-7.5 7.5M21 12H3"/></svg>
            E-commerce
          </div>
          <div class="port-title">Hi-Tec Bearings: 5,000+ SKUs, 300% organic traffic growth</div>
        </div>
      </div>
\end{lstlisting}

Line 293 carries \code{id="portfolio"}, which is what makes the hero's ``See the
work'' button work. An \code{id} is a unique name for one element, and a link to
\code{\#portfolio} scrolls to whatever carries that name.

\begin{jargon}{id, and how it differs from class}
An \code{id} names exactly one element on the page and must be unique. A
\code{class} labels any number of them. Use an \code{id} when something needs to
be linked to or found individually; use a \code{class} when you want to style a
group. This page has exactly two ids, \code{portfolio} and \code{book}, and both
exist to be jumped to.
\end{jargon}

Each card is built from two halves. The top half, \code{port-thumb}, is where a
real video thumbnail would go; today it is a gradient supplied by the class
\code{pt1}, defined at line 504 of the stylesheet. Over it sit a badge in the
corner and an overlay containing a play button, which the stylesheet keeps
invisible until the mouse is over the card. The bottom half, \code{port-info},
holds a small platform label with an arrow icon and the project title.

The six projects named across these cards are Hi-Tec Bearings, DigiMarvixx,
Synergy Event, IBA GPT, Prestige Sleep and IBA CED. Unlike the hero tiles, these
are specific, named and checkable, which is a meaningfully different kind of
claim.

\begin{honest}{The play buttons promise videos that do not exist}
Every card has a play button that appears on hover, and none of them plays
anything: the card is not a link and has no click handler. The stylesheet even
sets \code{cursor: pointer} on it (line 465), so the pointer changes to a hand,
promising an action that never happens. \file{ASSETS.md} confirms the intent was
real video, and the video was never supplied. As it stands the hover behaviour is
a small deception, and the cheapest honest fix is to remove the play overlay
until the videos exist.
\end{honest}

\paragraph{Lines 319 to 354: portfolio cards two and three.}

\begin{lstlisting}[language=HTML,firstnumber=319,caption={\file{index.html}, lines 319--354}]

      <div class="port-card" style="--i:1">
        <div class="port-thumb pt2">
          <div class="port-badge">3x leads</div>
          <div class="port-thumb-overlay">
            <div class="port-play">
              <svg viewBox="0 0 24 24" fill="white" aria-hidden="true"><path d="M5.25 5.653c0-.856.917-1.398 1.667-.986l11.54 6.348a1.125 1.125 0 010 1.971l-11.54 6.347a1.125 1.125 0 01-1.667-.985V5.653z"/></svg>
            </div>
          </div>
        </div>
        <div class="port-info">
          <div class="port-platform">
            <svg viewBox="0 0 24 24" fill="none" stroke="currentColor" stroke-width="1.8" aria-hidden="true"><path stroke-linecap="round" stroke-linejoin="round" d="M13.5 4.5L21 12m0 0l-7.5 7.5M21 12H3"/></svg>
            Web + WhatsApp
          </div>
          <div class="port-title">DigiMarvixx: 3x monthly WhatsApp leads, 60% mobile traffic</div>
        </div>
      </div>

      <div class="port-card" style="--i:2">
        <div class="port-thumb pt3">
          <div class="port-badge">1.2s load</div>
          <div class="port-thumb-overlay">
            <div class="port-play">
              <svg viewBox="0 0 24 24" fill="white" aria-hidden="true"><path d="M5.25 5.653c0-.856.917-1.398 1.667-.986l11.54 6.348a1.125 1.125 0 010 1.971l-11.54 6.347a1.125 1.125 0 01-1.667-.985V5.653z"/></svg>
            </div>
          </div>
        </div>
        <div class="port-info">
          <div class="port-platform">
            <svg viewBox="0 0 24 24" fill="none" stroke="currentColor" stroke-width="1.8" aria-hidden="true"><path stroke-linecap="round" stroke-linejoin="round" d="M13.5 4.5L21 12m0 0l-7.5 7.5M21 12H3"/></svg>
            Events
          </div>
          <div class="port-title">Synergy Event: 800 participants, 98/100 Lighthouse score</div>
        </div>
      </div>
\end{lstlisting}

Structurally identical to card one, differing only in the stagger index
(\code{-{}-i:1} and \code{-{}-i:2}), the gradient class (\code{pt2}, \code{pt3}), the
badge text, the platform label and the title. Note that the same arrow icon is
used as the ``platform'' icon on every card regardless of what the platform
actually is, so it is decorative rather than informative.

``98/100 Lighthouse score'' on line 352 is a specific, verifiable claim.
Lighthouse is Google's page-quality auditing tool, and 98 out of 100 is a strong
result.

\paragraph{Lines 355 to 412: portfolio cards four, five and six.}

\begin{lstlisting}[language=HTML,firstnumber=355,caption={\file{index.html}, lines 355--412}]

      <div class="port-card" style="--i:3">
        <div class="port-thumb pt4">
          <div class="port-badge">AI chat</div>
          <div class="port-thumb-overlay">
            <div class="port-play">
              <svg viewBox="0 0 24 24" fill="white" aria-hidden="true"><path d="M5.25 5.653c0-.856.917-1.398 1.667-.986l11.54 6.348a1.125 1.125 0 010 1.971l-11.54 6.347a1.125 1.125 0 01-1.667-.985V5.653z"/></svg>
            </div>
          </div>
        </div>
        <div class="port-info">
          <div class="port-platform">
            <svg viewBox="0 0 24 24" fill="none" stroke="currentColor" stroke-width="1.8" aria-hidden="true"><path stroke-linecap="round" stroke-linejoin="round" d="M13.5 4.5L21 12m0 0l-7.5 7.5M21 12H3"/></svg>
            AI Assistant
          </div>
          <div class="port-title">IBA GPT: AI assistant answering questions about IBA university</div>
        </div>
      </div>

      <div class="port-card" style="--i:4">
        <div class="port-thumb pt5">
          <div class="port-badge">24/7 sales</div>
          <div class="port-thumb-overlay">
            <div class="port-play">
              <svg viewBox="0 0 24 24" fill="white" aria-hidden="true"><path d="M5.25 5.653c0-.856.917-1.398 1.667-.986l11.54 6.348a1.125 1.125 0 010 1.971l-11.54 6.347a1.125 1.125 0 01-1.667-.985V5.653z"/></svg>
            </div>
          </div>
        </div>
        <div class="port-info">
          <div class="port-platform">
            <svg viewBox="0 0 24 24" fill="none" stroke="currentColor" stroke-width="1.8" aria-hidden="true"><path stroke-linecap="round" stroke-linejoin="round" d="M13.5 4.5L21 12m0 0l-7.5 7.5M21 12H3"/></svg>
            Health + E-commerce
          </div>
          <div class="port-title">Prestige Sleep: clinic booking and product shop in one site</div>
        </div>
      </div>

      <div class="port-card" style="--i:5">
        <div class="port-thumb pt6">
          <div class="port-badge">4 dashboards</div>
          <div class="port-thumb-overlay">
            <div class="port-play">
              <svg viewBox="0 0 24 24" fill="white" aria-hidden="true"><path d="M5.25 5.653c0-.856.917-1.398 1.667-.986l11.54 6.348a1.125 1.125 0 010 1.971l-11.54 6.347a1.125 1.125 0 01-1.667-.985V5.653z"/></svg>
            </div>
          </div>
        </div>
        <div class="port-info">
          <div class="port-platform">
            <svg viewBox="0 0 24 24" fill="none" stroke="currentColor" stroke-width="1.8" aria-hidden="true"><path stroke-linecap="round" stroke-linejoin="round" d="M13.5 4.5L21 12m0 0l-7.5 7.5M21 12H3"/></svg>
            Analytics
          </div>
          <div class="port-title">IBA CED: analytics platform tracking entrepreneurship outcomes</div>
        </div>
      </div>

    </div><!-- /portfolio-grid -->
  </div>
</section>
\end{lstlisting}

The final three cards complete the set of six, following the same template with
indices 3, 4 and 5 and gradients \code{pt4} through \code{pt6}. Two of the six
projects are for the same institution (IBA GPT and the analytics platform), which
is the only repetition among them.

Line 410 closes the grid with a naming comment, and lines 411 and 412 close the
container and the section.

\paragraph{Lines 413 to 472: the client proof band.}

\begin{lstlisting}[language=HTML,firstnumber=413,caption={\file{index.html}, lines 413--472}]


<!-- ══════════════════════════════════════════════
     CLIENT PROOF
     UPDATE: replace initials and numbers with real client data
══════════════════════════════════════════════ -->
<section class="proof-section" aria-label="Client results">
  <div class="container">
    <div class="reveal" style="margin-bottom:48px;">
      <p class="section-kicker">Client results</p>
      <h2 class="section-title">Numbers that<br>speak for themselves.</h2>
    </div>

    <div class="proof-grid stagger">

      <div class="proof-card" style="--i:0">
        <div class="proof-header">
          <div class="proof-avatar pa1">K</div>
          <div>
            <div class="proof-name">Hi-Tec Bearings</div>
            <div class="proof-company">Industrial E-commerce</div>
          </div>
        </div>
        <div class="proof-metric">
          <span class="proof-num">300%</span>
          <span class="proof-unit">organic traffic growth</span>
        </div>
      </div>

      <div class="proof-card" style="--i:1">
        <div class="proof-header">
          <div class="proof-avatar pa2">D</div>
          <div>
            <div class="proof-name">DigiMarvixx</div>
            <div class="proof-company">Digital Marketing</div>
          </div>
        </div>
        <div class="proof-metric">
          <span class="proof-num">3x</span>
          <span class="proof-unit">monthly WhatsApp leads</span>
        </div>
      </div>

      <div class="proof-card" style="--i:2">
        <div class="proof-header">
          <div class="proof-avatar pa3">S</div>
          <div>
            <div class="proof-name">Synergy Event</div>
            <div class="proof-company">Events & Conferences</div>
          </div>
        </div>
        <div class="proof-metric">
          <span class="proof-num">800</span>
          <span class="proof-unit">participants, 98/100 Lighthouse</span>
        </div>
      </div>

    </div>
  </div>
</section>
\end{lstlisting}

Three cards, each with a coloured circle bearing an initial, a client name, an
industry, and one headline number. The numbers repeat what the portfolio cards
already said about the same three clients, which is deliberate: repetition is how
a landing page makes a claim stick.

Note line 430: the avatar for Hi-Tec Bearings shows the letter \textbf{K}, not H.
That is almost certainly a leftover from the placeholder set in the hero (K, H, M,
A), and it is the clearest surviving trace of the \code{UPDATE} comment on line
417 having been only partly acted on.

Line 461 contains \code{Events \& Conferences}. The ampersand is written as a bare
character rather than as the escape \code{\&amp;}. Browsers forgive this, and it
is technically invalid: \code{\&} begins a character reference in HTML. It works
here only because what follows is a space, so no valid reference can be formed.
It is worth knowing as a habit to avoid rather than a bug to fix.

\begin{jargon}{Character reference, or entity}
A way of writing characters that would otherwise be interpreted as markup:
\code{\&amp;} for an ampersand, \code{\&lt;} for a less-than sign. They exist
because those characters have structural meaning in HTML.
\end{jargon}

\paragraph{Lines 473 to 549: the services grid.}

\begin{lstlisting}[language=HTML,firstnumber=473,caption={\file{index.html}, lines 473--495}]


<!-- ══════════════════════════════════════════════
     SERVICES
══════════════════════════════════════════════ -->
<section class="services-section" aria-label="Services">
  <div class="container">
    <div class="reveal" style="margin-bottom:48px;">
      <p class="section-kicker">What we do</p>
      <h2 class="section-title">Everything your content<br>operation needs.</h2>
    </div>

    <div class="services-grid stagger">

      <div class="service-item" style="--i:0">
        <div class="service-icon">
          <svg viewBox="0 0 24 24" fill="none" stroke="currentColor" stroke-width="1.8" aria-hidden="true">
            <path stroke-linecap="round" stroke-linejoin="round" d="M15.75 10.5l4.72-4.72a.75.75 0 011.28.53v11.38a.75.75 0 01-1.28.53l-4.72-4.72M4.5 18.75h9a2.25 2.25 0 002.25-2.25v-9a2.25 2.25 0 00-2.25-2.25h-9A2.25 2.25 0 002.25 7.5v9a2.25 2.25 0 002.25 2.25z"/>
          </svg>
        </div>
        <div class="service-name">AI Personal Brand</div>
        <span class="service-tag">Your voice, published</span>
      </div>
\end{lstlisting}

The services grid follows the established shape: a revealed heading pair, then a
staggered grid. Each of the six items is an icon box, a name, and a small
uppercase tag line. The first item's icon is, once again, the same video-camera
path from tile 1.

The remaining five items are structurally identical and are quoted in full below,
since the coverage rule admits no ``and so on''.

\begin{lstlisting}[language=HTML,firstnumber=496,caption={\file{index.html}, lines 496--549: services two to six}]

      <div class="service-item" style="--i:1">
        <div class="service-icon">
          <svg viewBox="0 0 24 24" fill="none" stroke="currentColor" stroke-width="1.8" aria-hidden="true">
            <path stroke-linecap="round" stroke-linejoin="round" d="M21 21l-5.197-5.197m0 0A7.5 7.5 0 105.196 5.196a7.5 7.5 0 0010.607 10.607z"/>
          </svg>
        </div>
        <div class="service-name">AI SEO & Content</div>
        <span class="service-tag">Rank for what they ask</span>
      </div>

      <div class="service-item" style="--i:2">
        <div class="service-icon">
          <svg viewBox="0 0 24 24" fill="none" stroke="currentColor" stroke-width="1.8" aria-hidden="true">
            <path stroke-linecap="round" stroke-linejoin="round" d="M8.625 12a.375.375 0 11-.75 0 .375.375 0 01.75 0zm0 0H8.25m4.125 0a.375.375 0 11-.75 0 .375.375 0 01.75 0zm0 0H12m4.125 0a.375.375 0 11-.75 0 .375.375 0 01.75 0zm0 0h-.375M21 12c0 4.556-4.03 8.25-9 8.25a9.764 9.764 0 01-2.555-.337A5.972 5.972 0 015.41 20.97a5.969 5.969 0 01-.474-.065 4.48 4.48 0 00.978-2.025c.09-.457-.133-.901-.467-1.226C3.93 16.178 3 14.189 3 12c0-4.556 4.03-8.25 9-8.25s9 3.694 9 8.25z"/>
          </svg>
        </div>
        <div class="service-name">AI Chat Agent</div>
        <span class="service-tag">Reply in seconds</span>
      </div>

      <div class="service-item" style="--i:3">
        <div class="service-icon">
          <svg viewBox="0 0 24 24" fill="none" stroke="currentColor" stroke-width="1.8" aria-hidden="true">
            <path stroke-linecap="round" stroke-linejoin="round" d="M2.25 6.75c0 8.284 6.716 15 15 15h2.25a2.25 2.25 0 002.25-2.25v-1.372c0-.516-.351-.966-.852-1.091l-4.423-1.106c-.44-.11-.902.055-1.173.417l-.97 1.293c-.282.376-.769.542-1.21.38a12.035 12.035 0 01-7.143-7.143c-.162-.441.004-.928.38-1.21l1.293-.97c.363-.271.527-.734.417-1.173L6.963 3.102a1.125 1.125 0 00-1.091-.852H4.5A2.25 2.25 0 002.25 4.5v2.25z"/>
          </svg>
        </div>
        <div class="service-name">AI Voice Agent</div>
        <span class="service-tag">Answer calls 24/7</span>
      </div>

      <div class="service-item" style="--i:4">
        <div class="service-icon">
          <svg viewBox="0 0 24 24" fill="none" stroke="currentColor" stroke-width="1.8" aria-hidden="true">
            <path stroke-linecap="round" stroke-linejoin="round" d="M21.75 6.75v10.5a2.25 2.25 0 01-2.25 2.25h-15a2.25 2.25 0 01-2.25-2.25V6.75m19.5 0A2.25 2.25 0 0019.5 4.5h-15a2.25 2.25 0 00-2.25 2.25m19.5 0v.243a2.25 2.25 0 01-1.07 1.916l-7.5 4.615a2.25 2.25 0 01-2.36 0L3.32 8.91a2.25 2.25 0 01-1.07-1.916V6.75"/>
          </svg>
        </div>
        <div class="service-name">AI Email Outreach</div>
        <span class="service-tag">Follow up until they respond</span>
      </div>

      <div class="service-item" style="--i:5">
        <div class="service-icon">
          <svg viewBox="0 0 24 24" fill="none" stroke="currentColor" stroke-width="1.8" aria-hidden="true">
            <path stroke-linecap="round" stroke-linejoin="round" d="M9.53 16.122a3 3 0 00-5.78 1.128 2.25 2.25 0 01-2.4 2.245 4.5 4.5 0 008.4-2.245c0-.399-.078-.78-.22-1.128zm0 0a15.998 15.998 0 003.388-1.62m-5.043-.025a15.994 15.994 0 011.622-3.395m3.42 3.42a15.995 15.995 0 004.764-4.648l3.876-5.814a1.151 1.151 0 00-1.597-1.597L14.146 6.32a15.996 15.996 0 00-4.649 4.763m3.42 3.42a6.776 6.776 0 00-3.42-3.42"/>
          </svg>
        </div>
        <div class="service-name">AI Ad Creatives</div>
        <span class="service-tag">On-brand at volume</span>
      </div>

    </div>
  </div>
</section>
\end{lstlisting}

The six services are AI Personal Brand, AI SEO and Content, AI Chat Agent, AI
Voice Agent, AI Email Outreach and AI Ad Creatives. Their icons in order are a
video camera, a magnifying glass, a speech bubble, a telephone handset, an
envelope and a paintbrush, each one an obvious visual match for its service, which
is why they need no labels of their own.

Line 503 contains another unescaped ampersand in ``AI SEO \& Content'', the same
point as line 461.

Note the grid gets its dividing lines by a trick rather than by borders: the
stylesheet gives the grid a 2-pixel gap with the border colour as the grid's own
background, then gives each cell the page background (lines 608 to 618). The gaps
show the layer beneath, producing hairlines with no border rules at all. It is
elegant, and it is the kind of thing that looks broken if you change one of the
two backgrounds without the other.

\paragraph{Lines 550 to 580: the booking panel.}

\begin{lstlisting}[language=HTML,firstnumber=550,caption={\file{index.html}, lines 550--580}]


<!-- ══════════════════════════════════════════════
     BOOK A CALL
     UPDATE: replace YOUR_CAL_LINK with Salman's Cal.com or Calendly URL.
     For inline Cal.com embed, swap the anchor tag for the Cal embed snippet.
══════════════════════════════════════════════ -->
<section class="book-section" id="book" aria-label="Book a call">
  <div class="container">
    <div class="book-wrap reveal">
      <p class="book-eyebrow">Ready to grow?</p>
      <h2 class="book-title">30 minutes.<br>No slides.</h2>
      <p class="book-sub">Pick a time. We talk about your content goals.</p>
      <div class="book-cta-row">
        <!-- UPDATE: replace href with your Cal.com or Calendly link -->
        <a href="https://cal.com/davonex/30min" target="_blank" rel="noopener" class="btn btn-primary" style="padding:16px 36px;font-size:17px;">
          Book a call
          <svg width="18" height="18" viewBox="0 0 24 24" fill="none" stroke="currentColor" stroke-width="2.5" aria-hidden="true">
            <path stroke-linecap="round" stroke-linejoin="round" d="M13.5 4.5L21 12m0 0l-7.5 7.5M21 12H3"/>
          </svg>
        </a>
      </div>
      <p class="book-detail">
        <svg viewBox="0 0 24 24" fill="none" stroke="currentColor" stroke-width="2" aria-hidden="true">
          <path stroke-linecap="round" stroke-linejoin="round" d="M12 6v6h4.5m4.5 0a9 9 0 11-18 0 9 9 0 0118 0z"/>
        </svg>
        Free · <span>30 minutes · Zoom</span>
      </p>
    </div>
  </div>
</section>
\end{lstlisting}

This is the destination of every ``Book a call'' button on the page, via
\code{id="book"} on line 557. It is the point of the whole document.

Line 559 combines two classes on one element, \code{book-wrap} for the panel's
appearance and \code{reveal} so the JavaScript fades it in. That is the only place
in the file where a layout class and a behaviour class share an element.

Line 565 is the actual booking link, and three of its attributes matter.
\code{target="\_blank"} opens it in a new tab, so the visitor does not lose the
page. \code{rel="noopener"} is a security measure that must accompany it: without
it, the newly opened page is handed a reference back to this one and can navigate
it somewhere else, a trick known as tab-nabbing. Its presence here is correct and
is the kind of detail that is usually forgotten.

\begin{honest}{The booking link is a placeholder that was never replaced}
\code{cal.com/davonex/30min} is the calendar of Davonex, a different agency
project, not a Digitalise AI calendar. Commit \code{4dffc12} on 2026-07-28 set
this deliberately, and the same placeholder calendar is recorded as an outstanding
item on the agency's main website too, so it is a known shared loose end rather
than a slip here.

Two things around it are worse than the link itself. The comment on line 554 still
tells the reader to search for \code{YOUR\_CAL\_LINK}, which no longer exists
anywhere in the file, and \file{ASSETS.md} repeats that same stale instruction.
Somebody following the written instructions would search, find nothing, and
conclude the work was done.
\end{honest}

Lines 572 to 577 are the reassurance line under the button: a clock icon, then
``Free'', then the duration and the platform. The middle dots here are the
U+00B7 character discussed in section 1.4, and they render exactly.

\paragraph{Lines 581 to 597: the footer, and the script tag.}

\begin{lstlisting}[language=HTML,firstnumber=581,caption={\file{index.html}, lines 581--597}]


<!-- ══════════════════════════════════════════════
     FOOTER
══════════════════════════════════════════════ -->
<footer>
  <div class="container">
    <p>
      <a href="/">digitalise.agency</a> &nbsp;·&nbsp;
      <a href="mailto:hello@digitalise.agency">hello@digitalise.agency</a>
    </p>
  </div>
</footer>

<script src="/js/main.js"></script>
</body>
</html>
\end{lstlisting}

The footer is one line of text: a link home and an email address, separated by a
middle dot padded with \code{\&nbsp;} on each side.

\begin{jargon}{Non-breaking space}
\code{\&nbsp;} is a space that a browser will not use as a place to wrap a line,
and which is not collapsed when several spaces sit together. Here it does the
second job: ordinary spaces around the dot would collapse to one, and these keep
the separator visually padded.
\end{jargon}

Line 589 links to \code{/}, meaning the site's own home page. The visible text
says \code{digitalise.agency}, which is a different site from
\code{ai.digitalise.agency}. So the footer link appears to promise the main agency
site and in fact reloads this one. On the live site that would be a small
misdirection; here it is another loose end.

Line 590 is a \code{mailto:} link, which opens the visitor's email program with
the address filled in.

Line 595 loads the JavaScript, and \textbf{where it sits is the whole point}. It
is the last thing before \code{</body>}, so by the time the browser reaches it,
every element on the page already exists. \file{main.js} begins by searching for
elements immediately, with no wrapper telling it to wait, and that only works
because of this placement. Move this one line up into the \code{<head>} and the
script would run against an empty document, find nothing, and the page would
quietly lose all three of its behaviours with no error message.

Lines 596 and 597 close the body and the document.

\subsubsection*{What it connects to}

\textbf{This file names:} \file{new-design-reference/css/site.css} (line 16),
\file{new-design-reference/js/main.js} (line 595), Google Fonts (lines 13 to 15),
and cal.com (line 565).

\textbf{This file is named by:} nothing inside the repository. The deploy workflow
copies the whole folder rather than referring to files individually, and
\file{netlify.toml} publishes the folder as a whole. \file{ASSETS.md} discusses it
at length but cannot affect it.

\textbf{Internal links:} the hero and navigation both point at \code{\#book}
(lines 33 and 162), and the hero also points at \code{\#portfolio} (line 168).
Those two ids, on lines 557 and 293, are the only link targets in the file.

\subsubsection*{How it breaks}

\begin{itemize}
\item \textbf{Rename a class here without renaming it in the stylesheet} and that
  element loses its appearance silently. Nothing warns you.
\item \textbf{Move the \code{<script>} tag into the \code{<head>}} and all three
  behaviours stop working, with no error in the console. The metrics band would
  then permanently display four zeros.
\item \textbf{Reorder the eight hero tiles} and they jump to different positions,
  because the stylesheet positions them by their number, not by their class.
\item \textbf{Open the file directly off a disk} and it renders unstyled, because
  of the leading slash on line 16. This looks like catastrophic breakage and is
  not. See section 11.
\item \textbf{Disable JavaScript} and the page is fully readable but static: the
  metrics show zeros, nothing fades in, and the anchor links jump instead of
  gliding. Only the metrics are a real loss.
\item \textbf{Delete the id on line 557} and every ``Book a call'' button on the
  page silently does nothing, because \file{main.js} returns early when it cannot
  find a link's target.
\end{itemize}

\subsubsection*{Honest findings for this file}

Collected here, each already argued above: the invented tile metrics (lines 55,
68, 74, 91, 116, 141); the missing \code{og:image} and favicon; the mismatched
avatar initial on line 429; five stale \code{UPDATE} comments (lines 187, 290,
417, 554, 564); the placeholder booking link on line 565; two unescaped ampersands
(lines 461 and 503); \code{<em>} used as a colour hook on line 158; play buttons
that play nothing; a footer link whose text and destination disagree; the
literal \code{\#34d399} used four times where a token exists; and one inline
padding that permanently overrides a stylesheet rule.
